\subsection{Steganalysis}
\label{app:subsec:steganalysis}
\paragraph{Raw Perplexity Results} \autoref{tab:perplexity-stego-c2} presents raw perplexity GPT-2 Large perplexity scores for stegotexts (\stegotext) and \typetwo covertexts (\covertext) for all three tasks. Broadly, \covertext and \stegotext perplexities are very similar across all three tasks, reinforcing the chance performance of the perplexity-based classifier (\S\ref{sec:steganalysis}).

\begin{table}[t]
\centering
\small
\setlength{\tabcolsep}{5pt}
\begin{tabular}{lcc}
\toprule
Text type & PPL (mean $\pm$ std) & Median \\
\midrule
\multicolumn{3}{l}{\textit{\qa}} \\
Stego    & $12.21 \pm 2.06$ & 11.94 \\
\typetwo & $10.70 \pm 1.77$ & 10.58 \\
\midrule
\multicolumn{3}{l}{\textit{\sg}} \\
Stego    & $18.69 \pm 2.84$ & 18.72 \\
\typetwo & $21.27 \pm 3.36$  & 20.93 \\
\midrule
\multicolumn{3}{l}{\textit{\lr}} \\
Stego    & $14.40 \pm 2.14$ & 14.42 \\
\typetwo & $15.15 \pm 2.33$ & 15.14 \\
\bottomrule
\end{tabular}
\caption{GPT-2 large perplexity, stego vs prompted length-matched \typetwo covertexts.}
\label{tab:perplexity-stego-c2}
\end{table}

\paragraph{AUC Results} \autoref{fig:steg-classifier-auc} Reports the area under the curve (AUC) for the same results as are plotted in \autoref{fig:steg-classifier-acc}.

\begin{figure}
    \centering
    \includegraphics[width=\linewidth]{figures/steg_auc_results.png}
    \caption{Perplexity (PPL) and Embedding-Based (Embed) classifier AUC in distinguishing stegotexts from \typeone (left) and \typetwo covertexts (right), based on the same results as \autoref{fig:steg-classifier-acc}}
    \label{fig:steg-classifier-auc}
\end{figure}

\paragraph{LLM as Judge}
An LLM is given a text and asked to decide if it is a stegotext that contains encrypted message or a generic LLM response. The trials are independent from each other and 150 examples from each type of texts are drawn from the dataset to be evaluated. Since using the same class of models as generation, namely GPT-4.1, would introduce self-recognition and common-prior leakage, where the model detects the local model's bias directly, we used Claude Sonnet 4.5 through OpenRouter for this experiment.

\begin{table}[t]
\centering
\small
\setlength{\tabcolsep}{5pt}
\begin{tabular}{lcccc}
\toprule
System & Acc. & \%High & Acc@High & Acc@Med \\
\midrule
\multicolumn{5}{l}{\textit{S vs \typeone}} \\
\qa   & 0.507 & 68.7\% & 0.505 & 0.511 \\
\sg & 0.517 & 39.7\% & 0.571 & 0.480 \\
\lr & 0.493 & 73.7\% & 0.502 & 0.474 \\
\midrule
\multicolumn{5}{l}{\textit{S vs \typetwo}} \\
\qa   & 0.560 & 84.7\% & 0.602 & 0.302 \\
% \sg & 0.777 & 47.7\% & 0.916 & 0.653 \\
\sg & 0.657 & 41.0\% & 0.780 & 0.586 \\
\lr & 0.527 & 86.0\% & 0.570 & 0.262 \\
% \midrule
% \multicolumn{5}{l}{\textit{S vs \typethree}} \\
% \sg & 0.657 & 41.0\% & 0.780 & 0.586 \\
\bottomrule
\end{tabular}
\caption{LLM-as-judge (Claude Sonnet 4.5) on the same 150-stego/150-cover
splits used by the trained classifiers. Each judgment includes a self-reported
confidence label; we report the share flagged \emph{high} (\%High) and
accuracy stratified by confidence (Acc@High / Acc@Med).}
\label{tab:steg-llm-judge}
\end{table}

 Against \typeone covertexts (same-pipeline, random-bit) every system is at chance (0.493--0.517), and the judge's high-confidence rate is 39--74\% while Acc@High barely exceeds Acc@Med, which means that model's self evaluated confidence is essentially uncalibrated when no real signal is present, and the slot-injection mechanism itself leaves no surface trace the judge can exploit. 
 
 Against \typetwo covers, \qa and \lr remain near chance whereas \sg jumps to 0.66, showing that the judge's discriminating signal is content-style rather than the pipeline fingerprint, and is only exposed when the stego pipeline's content channel commits to a generator distinct from the covertext's. Confidence is also better calibrated here, with Acc@High consistently above Acc@Med by 15--30 points across systems. 
 %The \typethree result shows that when the covertext is generated by the same upstream model as the stego content channel (Qwen outline + GPT-4.1 synthesis), \sg accuracy drops from 0.777 to 0.657, confirming that matching the ideation fingerprint from the local model is what restores the security against a strong LLM judge.

Overall, two of the attacks requires a dedicated dataset for training, and can not distinguish stegotexts generated through our pipeline. For \sg, it has been shown that the classifiers rely heavily on the local model's fingerprint, which could be easily removed by rotating the choice of local model.

\subsection{Attack Analysis}
\label{app:subsec:attack-analysis}
To evaluate the severity and other qualitative differences among three types of attacks. We used the attacked stegotexts from the main experiments as the samples and measured BERTScore F1 \cite{zhang2020bertscore}, BLEU \cite {papineni-etal-2002-bleu}, and TER (Translation Error Rate) to evaluate the attack impact on the stegotexts.

As shown in Table \ref{tab:attack-metrics}, BERTScore F1 exceeds 90\% for most attacks across systems across attacks, which suggests that the meaning of the stegotexts are largely preserved undergoing attacks. It also reveals that the global paraphrase by an LLM introduces large structural changes and meaning losses comparing with local paraphrase attack with $p$ set at 100\%. 

\syn, being the only attack that does not preserve semantic meaning, indeed yields the lowest BERTScores. In our  \rtt attack yields the highest BERTScores and BLEU scores, suggesting that it is the mildest attack of the three.

Global \llmp attack introduces structural changes to the stegotexts. For all three systems, even when global \llmp attack the stronger meaning loss for attacked stegotexts than local \llmp attacked ones with $p$ set to 100\%, the BLEU scores are still lower than the local \llmp attacked ones. Together with the higher TER scores, the results suggest that the discrepancies between the stegotexts and globally paraphrased texts is more structural than lexical.

\begin{table}
\centering
\small
\setlength{\tabcolsep}{5pt}
\renewcommand{\arraystretch}{1.05}
\begin{tabular}{llccc}
\toprule
Attack & $p$ & BERTScore F1 & BLEU & TER \\
\midrule
\multicolumn{5}{l}{\textbf{Long-From Q\&A}} \\
Synonym                 & 0.2 & 0.951 & 71.2 & 13.2 \\
Synonym                 & 0.5 & 0.893 & 36.5 & 34.2 \\
Synonym                 & 1.0 & 0.802 &  4.2 & 68.0 \\
Local paraphrase        & 0.2 & 0.983 & 82.8 & 16.2 \\
Local paraphrase        & 0.5 & 0.958 & 55.4 & 42.8 \\
Local paraphrase        & 1.0 & 0.922 &  7.0 & 86.6 \\
Local translation  & 0.2 & 0.992 & 91.8 &  5.4 \\
Local translation  & 0.5 & 0.982 & 79.5 & 13.9 \\
Local translation  & 1.0 & 0.967 & 60.5 & 27.2 \\
Global paraphrase       & 1.0 & 0.905 & 11.5 & 88.3 \\
Global translation & 1.0 & 0.977 & 66.4 & 22.9 \\
\midrule
\multicolumn{5}{l}{\textbf{Story Generation}} \\
Synonym                 & 0.2 & 0.948 & 73.0 & 12.8 \\
Synonym                 & 0.5 & 0.894 & 40.1 & 32.9 \\
Synonym                 & 1.0 & 0.810 &  6.4 & 64.8 \\
Local paraphrase        & 0.2 & 0.979 & 82.2 & 16.4 \\
Local paraphrase        & 0.5 & 0.949 & 55.2 & 43.6 \\
Local paraphrase        & 1.0 & 0.910 &  9.9 & 87.1 \\
Local translation       & 0.2 & 0.986 & 89.5 & 7.32 \\
Local translation       & 0.5 & 0.971 & 74.4 & 18.4 \\
Local translation       & 1.0 & 0.950 & 51.5 & 35.7\\
Global paraphrase       & 1.0 & 0.888 & 18.7 & 88.9 \\
Global translation      & 1.0 & 0.967 & 63.0 & 26.4 \\
\midrule
\multicolumn{5}{l}{\textbf{Literature Review}} \\
Synonym                 & 0.2 & 0.960 & 75.8 & 11.7 \\
Synonym                 & 0.5 & 0.918 & 47.4 & 28.9 \\
Synonym                 & 1.0 & 0.840 & 16.9 & 57.9 \\
Local paraphrase        & 0.2 & 0.980 & 82.3 & 17.1 \\
Local paraphrase        & 0.5 & 0.960 & 61.8 & 37.7 \\
Local paraphrase        & 1.0 & 0.923 & 19.3 & 81.2 \\
Local translation       & 0.2 & 0.991 & 92.6 &  5.5 \\
Local translation       & 0.5 & 0.979 & 81.1 & 14.4 \\
Local translation       & 1.0 & 0.962 & 63.9 & 28.1 \\
Global paraphrase       & 1.0 & 0.912 & 33.8 & 83.7 \\
Global translation      & 1.0 & 0.982 & 77.6 & 15.9 \\
\bottomrule
\end{tabular}
\caption{Attack severity metrics on a subset of the steganalysis dataset
(\qa 6\,bits, \sg 18\,bits, \lr 20\,bits). Each row aggregates 30 stegotexts with 3 runs per stochastic attack ($n{=}90$); only one run ($n{=}30$) is used for \syn attack as it is deterministic within runs.}
\label{tab:attack-metrics}
\end{table}
